% GENERATED FILE. Edit canonical lessons, then run npm run generate:books.
% canonical-source-hash: fnv1a64:f1acc3c486caf278
% canonical-lessons: FA-C529-nevisande, FA-C529-montaqed, FA-C529-naqd-kardan, FA-C529-motalee-kardan, FA-C529-kerm-e-ketab

\chapter{Readers and Critics}
\label{ch:fa-readers-and-critics}
\hlchaptermodality{529}

\begin{chapteropening}
\textbf{By the end of this chapter:} \emph{I can talk about writers, readers and reviews.}

Complete the last lesson of chapter 529: I can talk about writers, readers and reviews.
\end{chapteropening}

\section[\texorpdfstring{\fa{نویسنده} (nevisande)}{نویسنده (nevisande)}]{\fa{نویسنده} (nevisande) — a writer, an author}

\label{lesson:FA-C529-nevisande}



\begin{quote}
Before the new one: say the Persian for a publisher, then the Persian for a translator.
\end{quote}

\subsection*{You'll want to know: \fa{نویسنده}}
\textbf{\fa{نویسنده}} (\emph{nevisande}) — \textquotedblleft{}a writer, an author\textquotedblright{}.

\begin{quote}
\fa{این} \fa{نویسنده} \fa{برای} \fa{بچه} \fa{ها} \fa{داستان} \fa{مینویسد}.
\end{quote}

\begin{quote}
\emph{in nevisande barâ-ye bachchehâ dâstân minevisad.}
\end{quote}

\begin{quote}
This writer writes stories for children.
\end{quote}

\emph{nevis-} (stem of \emph{neveshtan}) + \emph{-ande} 'one who'.

\begin{cousinweb}
\emph{nevisande} is built on native Persian \emph{nevis-}, the present stem of \emph{neveshtan}, 'to write', which EtymDB traces to Middle Persian (EtymDB), with the ending \emph{-ande}, 'one who'.
\end{cousinweb}

\begin{grammarlens}[title={What you've built}]
One more word for telling, explaining and describing at B1.
\end{grammarlens}

\subsection*{Your turn}
\begin{itemize}
\raggedright
  \item \emph{Say it:} \emph{nevisande}
  \item \emph{Say it:} \emph{nevisande}, once more
  \item \emph{Say it:} a sentence of your own with \emph{nevisande}
  \item \emph{Recall:} say the Persian for a publisher, then the Persian for a translator, then say \emph{nevisande} again
\end{itemize}

\begin{tcolorbox}[breakable,colback=teal!4,colframe=teal!35!black,title={Before you move on}]
What does \fa{نویسنده} mean? (\textquotedblleft{}A writer, an author\textquotedblright{}.) Say it once more.
\end{tcolorbox}
\section[\texorpdfstring{\fa{منتقد} (montaqed)}{منتقد (montaqed)}]{\fa{منتقد} (montaqed) — a critic}

\label{lesson:FA-C529-montaqed}



\begin{quote}
Before the new one: say the Persian for a translator, then the Persian for a writer.
\end{quote}

\subsection*{You'll want to know: \fa{منتقد}}
\textbf{\fa{منتقد}} (\emph{montaqed}) — \textquotedblleft{}a critic\textquotedblright{}.

\begin{quote}
\fa{منتقدها} \fa{از} \fa{فیلم} \fa{تازه} \fa{او} \fa{تعریف} \fa{کردند}.
\end{quote}

\begin{quote}
\emph{montaqedhâ az film-e tâze-ye u ta'rif kardand.}
\end{quote}

\begin{quote}
Critics praised his new film.
\end{quote}

Plural \emph{montaqedân} or \emph{montaqedhâ}. \emph{naqd} is criticism.

\begin{cousinweb}
\emph{montaqed} comes from Arabic \emph{muntaqid} (EtymDB), a word from the same root as \emph{naqd}, 'criticism'.
\end{cousinweb}

\begin{grammarlens}[title={What you've built}]
One more word for telling, explaining and describing at B1.
\end{grammarlens}

\subsection*{Your turn}
\begin{itemize}
\raggedright
  \item \emph{Say it:} \emph{montaqed}
  \item \emph{Say it:} \emph{montaqed}, once more
  \item \emph{Say it:} a sentence of your own with \emph{montaqed}
  \item \emph{Recall:} say the Persian for a translator, then the Persian for a writer, then say \emph{montaqed} again
\end{itemize}

\begin{tcolorbox}[breakable,colback=teal!4,colframe=teal!35!black,title={Before you move on}]
What does \fa{منتقد} mean? (\textquotedblleft{}A critic\textquotedblright{}.) Say it once more.
\end{tcolorbox}
\section[\texorpdfstring{\fa{نقد} \fa{کردن} (naqd kardan)}{نقد کردن (naqd kardan)}]{\fa{نقد} \fa{کردن} (naqd kardan) — to review, to critique}

\label{lesson:FA-C529-naqd-kardan}



\begin{quote}
Before the new one: say the Persian for a writer, then the Persian for a critic.
\end{quote}

\subsection*{You'll want to know: \fa{نقد} \fa{کردن}}
\textbf{\fa{نقد} \fa{کردن}} (\emph{naqd kardan}) — \textquotedblleft{}to review, to critique\textquotedblright{}.

\begin{quote}
\fa{او} \fa{فیلمها} \fa{را} \fa{در} \fa{مجله} \fa{نقد} \fa{میکند}.
\end{quote}

\begin{quote}
\emph{u filmhâ râ dar majalle naqd mikonad.}
\end{quote}

\begin{quote}
He reviews films in a magazine.
\end{quote}

\emph{naqd} 'criticism, review' + \emph{kardan}. The same word \emph{naqd} also means 'cash'.

\begin{cousinweb}
\emph{naqd} is a loan, usually traced to Arabic \emph{naqd}, 'ready money; testing coins'; testing a coin is usually taken as the picture behind testing a book or a film. EtymDB has no entry for it. \emph{kardan} is inherited Persian.
\end{cousinweb}

\begin{grammarlens}[title={What you've built}]
One more word for telling, explaining and describing at B1.
\end{grammarlens}

\subsection*{Your turn}
\begin{itemize}
\raggedright
  \item \emph{Say it:} \emph{naqd kardan}
  \item \emph{Say it:} \emph{naqd kardan}, once more
  \item \emph{Say it:} a sentence of your own with \emph{naqd kardan}
  \item \emph{Recall:} say the Persian for a writer, then the Persian for a critic, then say \emph{naqd kardan} again
\end{itemize}

\begin{tcolorbox}[breakable,colback=teal!4,colframe=teal!35!black,title={Before you move on}]
What does \fa{نقد} \fa{کردن} mean? (\textquotedblleft{}To review, to critique\textquotedblright{}.) Say it once more.
\end{tcolorbox}
\section[\texorpdfstring{\fa{مطالعه} \fa{کردن} (motâle'e kardan)}{مطالعه کردن (motâle'e kardan)}]{\fa{مطالعه} \fa{کردن} (motâle'e kardan) — to study (closely), to read}

\label{lesson:FA-C529-motalee-kardan}



\begin{quote}
Before the new one: say the Persian for a critic, then the Persian for to review.
\end{quote}

\subsection*{You'll want to know: \fa{مطالعه} \fa{کردن}}
\textbf{\fa{مطالعه} \fa{کردن}} (\emph{motâle'e kardan}) — \textquotedblleft{}to study (closely), to read\textquotedblright{}.

\begin{quote}
\fa{او} \fa{هر} \fa{شب} \fa{قبل} \fa{از} \fa{خواب} \fa{مطالعه} \fa{میکند}.
\end{quote}

\begin{quote}
\emph{u har shab qabl az khâb motâle'e mikonad.}
\end{quote}

\begin{quote}
She reads every night before bed.
\end{quote}

\emph{motâle'e} 'study, reading' + \emph{kardan}. More serious than \emph{khândan}.

\begin{cousinweb}
\emph{motâle'e} is a loan, usually traced to Arabic \emph{muṭālaʿa}, 'perusal, close reading'. EtymDB has no entry for it. \emph{kardan} is inherited Persian.
\end{cousinweb}

\begin{grammarlens}[title={What you've built}]
One more word for telling, explaining and describing at B1.
\end{grammarlens}

\subsection*{Your turn}
\begin{itemize}
\raggedright
  \item \emph{Say it:} \emph{motâle'e kardan}
  \item \emph{Say it:} \emph{motâle'e kardan}, once more
  \item \emph{Say it:} a sentence of your own with \emph{motâle'e kardan}
  \item \emph{Recall:} say the Persian for a critic, then the Persian for to review, then say \emph{motâle'e kardan} again
\end{itemize}

\begin{tcolorbox}[breakable,colback=teal!4,colframe=teal!35!black,title={Before you move on}]
What does \fa{مطالعه} \fa{کردن} mean? (\textquotedblleft{}To study (closely), to read\textquotedblright{}.) Say it once more.
\end{tcolorbox}
\section[\texorpdfstring{\fa{کرم} \fa{کتاب} (kerm-e ketâb)}{کرم کتاب (kerm-e ketâb)}]{\fa{کرم} \fa{کتاب} (kerm-e ketâb) — a bookworm}

\label{lesson:FA-C529-kerm-e-ketab}



\begin{quote}
Before the new one: say the Persian for to review, then the Persian for to study.
\end{quote}

\subsection*{You'll want to know: \fa{کرم} \fa{کتاب}}
\textbf{\fa{کرم} \fa{کتاب}} (\emph{kerm-e ketâb}) — \textquotedblleft{}a bookworm\textquotedblright{}.

\begin{quote}
\fa{خواهرم} \fa{کرم} \fa{کتاب} \fa{است} \fa{و} \fa{هر} \fa{هفته} \fa{یک} \fa{رمان} \fa{میخواند}.
\end{quote}

\begin{quote}
\emph{khâharam kerm-e ketâb ast va har hafte yek român mikhânad.}
\end{quote}

\begin{quote}
My sister is a bookworm and reads a novel every week.
\end{quote}

\emph{kerm} 'worm' + \emph{ketâb} 'book', as in English.

\begin{cousinweb}
\emph{kerm-e ketâb} joins inherited \emph{kerm}, 'worm', which EtymDB traces to Proto-Iranian and sets beside Lithuanian \emph{kirmis}, 'worm' (EtymDB), and \emph{ketâb}, 'book', from Arabic.
\end{cousinweb}

\begin{grammarlens}[title={What you've built}]
One more word for telling, explaining and describing at B1.
\end{grammarlens}

\subsection*{Your turn}
\begin{itemize}
\raggedright
  \item \emph{Say it:} \emph{kerm-e ketâb}
  \item \emph{Say it:} \emph{kerm-e ketâb}, once more
  \item \emph{Say it:} a sentence of your own with \emph{kerm-e ketâb}
  \item \emph{Recall:} say the Persian for to review, then the Persian for to study, then say \emph{kerm-e ketâb} again
\end{itemize}

\begin{tcolorbox}[breakable,colback=teal!4,colframe=teal!35!black,title={Before you move on}]
What does \fa{کرم} \fa{کتاب} mean? (\textquotedblleft{}A bookworm\textquotedblright{}.) Say it once more.
\end{tcolorbox}
